\section{A positive-Hilbert origin of the balanced determinant spectrum}
\label{sec:positive-hilbert-balance}

\textbf{Status: \StatusModel\ for a balanced positive-Hilbert multiplet and
its exact or approximate selective Ward protection over an arbitrary
positive interacting operator on every member of a supplied finite geometry
language, including a positive finite geometry channel; \StatusOpen\ for the
microscopic origin of the protected commutant, a dynamically selected
geometry language and projective refinement measure, continuum Lorentz
spin--statistics and anomaly freedom, nonlinear BV closure, and a
background-independent continuum quantum measure.}

The regulator theorem in
Theorem~\ref{thm:three-newton-universality} identifies exactly which
spectral moments must vanish, but it does not explain why a physical
microscopic system should have the signed weights needed for that
cancellation.  Entering \((1,-2,1)\) as a list of ``multiplicities'' is not
an answer: a negative multiplicity is not a positive state space.  The
companion construction \cite{XuPositiveHilbertBalance2026} closes this
narrower free-model problem by separating two notions that must not be
conflated:
\[
 \text{sign of a closed Gaussian loop}
 \quad\ne\quad
 \text{sign of a state norm}.
\]

\subsection{Exterior parity and an additive mass ladder}

Let \(G\) be any finite weighted graph and let \(K_G\geq0\) be one real
symmetric local kernel on it.  Introduce the positive internal Hilbert space
\begin{equation}
 \Hh_{\rm int}=\Lambda^\bullet\mathbb C^k
 =\bigoplus_{n=0}^{k}\Lambda^n\mathbb C^k,
 \qquad
 N=\sum_{a=1}^{k}c_a^\dagger c_a ,
 \label{eq:ph-exterior}
\end{equation}
with its standard exterior-product inner product.  The degree-\(n\)
subspace has dimension \(\binom{k}{n}\) and parity
\((-1)^F=(-1)^N\).  Give every component the same graph operator and the
additive mass-squared operator
\begin{equation}
 M^2=x\mathbf 1+\Delta N,\qquad x,\Delta>0,\qquad
 h_n=K_G+(x+n\Delta)\mathbf 1 .
 \label{eq:ph-mass}
\end{equation}
Even components are represented by commuting Gaussian variables and odd
components by Grassmann variables.  Integrating out the complete multiplet
therefore gives
\begin{equation}
 \Gamma_k[K_G]
 =\sum_{n=0}^{k}(-1)^n\binom{k}{n}
   \Tr\log\!\left[K_G+(x+n\Delta)\mathbf 1\right].
 \label{eq:ph-gamma}
\end{equation}
Both the multiplicity and the loop sign in this expression are fixed by
\(\Lambda^\bullet\mathbb C^k\); they are not independently tuned
coefficients.

\begin{theorem}[Exterior-algebra spectral balance]
\label{thm:ph-balance}
For Eqs.~\eqref{eq:ph-exterior} and \eqref{eq:ph-mass},
\begin{equation}
 \Str_{\Hh_{\rm int}}e^{-sM^2}
 =e^{-sx}(1-e^{-s\Delta})^k .
 \label{eq:ph-heat}
\end{equation}
Consequently,
\begin{equation}
 \Str(M^2)^j=0\quad(0\leq j<k),\qquad
 \Str(M^2)^k=(-1)^k k!\Delta^k .
 \label{eq:ph-moments}
\end{equation}
The identities hold at every graph size and for every common kernel
\(K_G\).
\end{theorem}

\begin{proof}
Resolving the supertrace by exterior degree gives
\[
 \Str e^{-sM^2}
 =e^{-sx}\sum_{n=0}^{k}(-1)^n\binom{k}{n}e^{-sn\Delta}.
\]
The binomial theorem yields Eq.~\eqref{eq:ph-heat}.  Since
\(1-e^{-s\Delta}=s\Delta+O(s^2)\), its zero at \(s=0\) has order exactly
\(k\).  Taking the first \(k\) derivatives at the origin proves
Eq.~\eqref{eq:ph-moments}.
\end{proof}

For the induced Newton term the minimal choice is \(k=2\):
\begin{equation}
 (w_0,w_1,w_2)=(1,-2,1),\qquad
 m_n^2=x+n\Delta,\qquad
 \Str\mathbf 1=\Str M^2=0 .
 \label{eq:ph-k2}
\end{equation}
The logarithmic invariant that survives the two cancellations is
\begin{align}
 F_2(x,\Delta)
 &:=\sum_{n=0}^{2}w_n m_n^2
          \log\frac{m_n^2}{\mu^2}\nonumber\\
 &=\int_0^\Delta\!\diff u\int_0^\Delta\!\diff v\,
     \frac{1}{x+u+v}>0 .
 \label{eq:ph-positive-invariant}
\end{align}
The integral is the second finite-difference identity for
\(f(y)=y\log(y/\mu^2)\).  It proves both positivity and independence of
\(\mu\), without subtracting large cutoff terms.  At \(x=\Delta=M^2\),
\begin{equation}
 F_2=M^2(3\log3-4\log2)
 =0.523248143764548\ldots\,M^2 .
 \label{eq:ph-value}
\end{equation}

\subsection{Where positivity is proved}

The determinant supertrace is not used as the physical inner product.
Associate a bosonic Fock factor \(\F_+(h_n)\) with each even component and a
fermionic Fock factor \(\F_-(h_n)\) with each odd component:
\begin{equation}
 \F_{\rm phys}
 =\bigotimes_{n\ {\rm even}}
   \F_+(h_n)^{\otimes\binom{k}{n}}
 \otimes
 \bigotimes_{n\ {\rm odd}}
   \F_-(h_n)^{\otimes\binom{k}{n}} .
 \label{eq:ph-fock}
\end{equation}
Every factor has its standard positive norm.  On the finite graph,
\begin{equation}
 \widehat H=\sum_{n,\alpha}
 \diff\Gamma_\pm(h_n)_{n,\alpha}\geq0,\qquad
 T_a=e^{-a\widehat H},\quad 0<T_a\leq\mathbf1 .
 \label{eq:ph-transfer}
\end{equation}

\begin{proposition}[Finite transfer-form positivity]
\label{prop:ph-positive}
For every positive density operator \(\rho\) and finite operator family
\(\{A_\alpha\}\), the matrix
\begin{equation}
 (G_T)_{\alpha\beta}
 =\Tr\!\left[
 (T_a^{1/2}A_\alpha\rho^{1/2})^\dagger
 (T_a^{1/2}A_\beta\rho^{1/2})\right]
 \label{eq:ph-gram}
\end{equation}
is positive semidefinite.  This statement is independent of the loop
weights \(w_n\).
\end{proposition}

\begin{proof}
For every vector \(z\),
\[
 z^\dagger G_Tz
 =\left\|T_a^{1/2}
 \left(\sum_\alpha z_\alpha A_\alpha\right)\rho^{1/2}\right\|_2^2
 \geq0.
\]
The spectral theorem gives \(\widehat H\geq0\) and \(0<T_a\leq\mathbf1\).
\end{proof}

Thus the minus signs in Eq.~\eqref{eq:ph-gamma} are ordinary closed
fermion-loop signs.  Proposition~\ref{prop:ph-positive} is a finite-graph
transfer statement, not an interacting Osterwalder--Schrader
reconstruction \cite{OsterwalderSchrader1973}.

\subsection{Regulator and same-link interfaces}

In the regulator class used in
Theorem~\ref{thm:three-newton-universality}, every scheme-dependent term is
proportional either to \(\Str\mathbf1\) or to \(\Str M^2\).  Hence the
two-mode multiplet gives the common limit
\begin{equation}
 Z_*=\kappa_g F_2(x,\Delta)>0 .
 \label{eq:ph-zstar}
\end{equation}
For \(x=\Delta=1\) and \(\Lambda=1200\), the archived sharp-momentum,
sharp-proper-time, and smooth-proper-time evaluations are respectively
\[
 0.5232467549,\qquad 0.5232474487,\qquad 0.5232388701,
\]
with a total scheme spread \(8.58\times10^{-6}\).  The exact target is
Eq.~\eqref{eq:ph-value}; these values are deterministic quadratures rather
than fits.

The same graph links can be varied rather than merely used as a common
spectral label.  If \(K_G(\eta)\) is differentiable, then
\begin{equation}
 \frac{\diff\Gamma_2}{\diff\eta}
 =\sum_{n=0}^{2}w_n
 \Tr\!\left[
 (K_G+m_n^2\mathbf1)^{-1}\frac{\diff K_G}{\diff\eta}\right].
 \label{eq:ph-same-link}
\end{equation}
On the archived inhomogeneous \(17\)-cycle, the analytic value is
\(0.022621032016492876\).  A central difference converges with measured
exponent \(1.99996\) and reaches absolute error
\(9.77\times10^{-11}\) before floating-point saturation.  This verifies
the weight--mass--link interface that is needed when the multiplet is
placed on the record support of Sec.~\ref{sec:same-link-induced-gravity}.

\subsection{Finite geometry sums and why exact pairing is insufficient}

Because the graph operator and internal mass act on different tensor
factors, the balance is pointwise in geometry.  Let \(\Omega\) be any finite
family of positive graph kernels \(K_\omega\), of possibly different
dimensions, with a supplied positive prior \(p_\omega\).  Define
\begin{equation}
 \Gamma_\omega
 =\sum_{n=0}^{k}w_n\Tr_\omega\log(K_\omega+m_n^2\mathbf1),
 \qquad
 \mathcal Z_\Omega
 =\sum_{\omega\in\Omega}p_\omega e^{-\Gamma_\omega}.
 \label{eq:ph-graph-sum}
\end{equation}

\begin{proposition}[Geometry-by-geometry spectral balance]
\label{prop:ph-graph-sum}
For every \(\omega\) and \(s>0\),
\begin{align}
 &\sum_{n=0}^{k}w_n
 \Tr_\omega e^{-s(K_\omega+m_n^2\mathbf1)}
 \nonumber\\
 &\hspace{1.5cm}
 =\Tr_\omega e^{-sK_\omega}\,
 e^{-sx}(1-e^{-s\Delta})^k .
 \label{eq:ph-graph-heat}
\end{align}
Thus the first \(k\) internal ultraviolet moments vanish on each graph
before geometry averaging, and \(\mathcal Z_\Omega\) is finite and strictly
positive.  For every finite move
\(K(t)=K_0+t\,\delta K\geq0\),
\begin{align}
 \Gamma[K_0+\delta K]-\Gamma[K_0]
 &=
 \int_0^1\!\diff t\sum_n w_n
 \nonumber\\[-2mm]
 &\quad\times
 \Tr[(K(t)+m_n^2\mathbf1)^{-1}\delta K] .
 \label{eq:ph-finite-move}
\end{align}
\end{proposition}

\begin{proof}
The heat operator factorizes on graph and internal tensor factors, and
Theorem~\ref{thm:ph-balance} evaluates the internal supertrace.  Every
shifted determinant is strictly positive, so every term in
Eq.~\eqref{eq:ph-graph-sum} has positive real weight.  The final identity
is the fundamental theorem of calculus applied to the log determinant.
\end{proof}

This extension is nonperturbative in the amplitude of a finite graph move
and uses no cancellation between geometries.  It is not yet a dynamical
sum over causal complexes: \(\Omega\), \(p_\omega\), and the refinement maps
are inputs.  In the supporting data, 64 connected weighted graphs with
7--15 vertices give a maximum pointwise heat-factorization residual
\(1.07\times10^{-14}\); an order-one chord insertion agrees with
Eq.~\eqref{eq:ph-finite-move} to \(2.98\times10^{-15}\).

The most obvious interaction-protection proposal also has a precise
failure mode.  Exact supersymmetry can be realized in interacting lattice
models and on arbitrary triangulations
\cite{CatterallKaramov2002,BerensteinCatterall2025}, but exact all-level
pairing cancels too much.

\begin{theorem}[Exact-pairing obstruction]
\label{thm:ph-pairing-obstruction}
Let \(\mathcal V=\mathcal V_+\oplus\mathcal V_-\) be a finite positive
graded Hilbert space, let \(D=D^\dagger\) be odd, and set
\(\mathcal H_D=D^2\).  Then for every continuous spectral function \(f\),
\begin{equation}
 \Str f(\mathcal H_D)=\operatorname{ind}(D_+)\,f(0).
 \label{eq:ph-pairing-index}
\end{equation}
In particular,
\(\Str[\mathcal H_D\log(\mathcal H_D/\mu^2)]=0\), with
\(0\log0:=0\).
\end{theorem}

\begin{proof}
On each positive \(\lambda\)-eigenspace,
\(D/\sqrt{\lambda}\) is an isometric bijection between even and odd
states, so their contributions cancel.  The zero-mode difference is
\(\operatorname{ind}(D_+)\).  The logarithmic function vanishes at zero.
\end{proof}

The required interacting symmetry is therefore selective: it must protect
\(\Str\mathbf1=\Str M^2=0\) but leave positive-energy spectra sufficiently
misaligned that \(F_2>0\).  Broken-supersymmetry mass-sum rules provide
continuum examples of partial rather than exact pairing
\cite{ItoyamaMaru2015,BenaEtAl2016}; they do not yet derive the required
record-lattice Ward identity.

\subsection{Selective Ward closure on a supplied finite geometry language}
\label{sec:selective-ward-closure}

The companion selective-Ward construction
\cite{XuSelectiveLadderWard2026} closes the finite interacting part of the
preceding problem without imposing exact spectral pairing.  Let
\(\K\) be any finite positive Hilbert space, let
\(\F_k=\Lambda^\bullet\mathbb C^k\), and let \(c_a^\dagger\), \(N\), and
\(\Pi=(-1)^N\) be the canonical creation, number, and parity operators on
\(\F_k\).  All physical interactions and all degrees of freedom of one
fixed geometry are contained in an arbitrary self-adjoint operator on
\(\K\).

\begin{theorem}[Selective ladder characterization and stability]
\label{thm:selective-ladder-closure}
For a self-adjoint \(H\) on \(\K\otimes\F_k\) and a real
\(\Delta\), the identities
\begin{equation}
 [H,\mathbf1\otimes c_a^\dagger]
 =\Delta(\mathbf1\otimes c_a^\dagger),
 \qquad a=1,\ldots,k,
 \label{eq:selective-ladder-ward}
\end{equation}
hold if and only if
\begin{equation}
 H=A\otimes\mathbf1+\Delta\,\mathbf1\otimes N
 \label{eq:selective-ladder-normal}
\end{equation}
for one self-adjoint \(A\) on \(\K\).  Consequently,
\begin{equation}
 \Str e^{-sH}=\Tr_\K e^{-sA}(1-e^{-s\Delta})^k,
 \qquad
 \Str H^j=0\quad(0\leq j<k).
 \label{eq:selective-heat}
\end{equation}
For an arbitrary self-adjoint \(H\), define
\[
 \varepsilon=\max_a
 \|[H,c_a^\dagger]-\Delta c_a^\dagger\|.
\]
Its Hilbert--Schmidt projection \(H_0\) onto the affine manifold
in Eq.~\eqref{eq:selective-ladder-normal} satisfies
\begin{equation}
 \|H-H_0\|\leq4k\varepsilon.
 \label{eq:selective-distance}
\end{equation}
If \(H,H_0\geq m\mathbf1\), \(D=2^k\dim\K\), and \(s>0\), then
\begin{equation}
 \left|\Str(e^{-sH}-e^{-sH_0})\right|
 \leq4kD\,s e^{-sm}\varepsilon .
 \label{eq:selective-heat-bound}
\end{equation}
\end{theorem}

\begin{proof}
After subtracting \(\Delta N\), Eq.~\eqref{eq:selective-ladder-ward}
commutes with the irreducible CAR algebra and therefore acts trivially on
\(\F_k\), proving the characterization.  Equation
\eqref{eq:selective-heat} follows by tensor factorization.  For the
approximate statement, average \(H-\Delta N\) over the finite Majorana
conjugation group.  The average is the conditional expectation onto the
CAR commutant, and a telescoping product of at most \(2k\) Majoranas bounds
the distance by \(4k\varepsilon\).  The Duhamel formula and
\(\|e^{-tH}\|\leq e^{-tm}\) give
Eq.~\eqref{eq:selective-heat-bound}.
\end{proof}

This theorem is nonperturbative in \(A\): \(A\) may contain arbitrary
finite interacting physical and geometric terms.  For \(k=2\), the finite
response is
\begin{align}
 \I_\Delta(A)
 &:=
 \Tr\!\left[
 A\log A-2(A+\Delta)\log(A+\Delta)\right.\nonumber\\[-1mm]
 &\hspace{29mm}\left.
 +(A+2\Delta)\log(A+2\Delta)\right]\nonumber\\
 &=\int_0^\Delta\!\diff u\int_0^\Delta\!\diff v\,
 \Tr(A+u+v)^{-1}>0 .
 \label{eq:selective-invariant}
\end{align}
Thus the protected low moments coexist with a nonzero response.  A fresh-cell
unitary of the factorized form
\[
 U_h=U_{\rm phys}\otimes e^{-ih\Delta N}
\]
preserves the shifted ladder eigenoperator exactly even when
\(U_{\rm phys}\) is strong.  This realizes the Ward law microscopically, but
also exposes its present boundary: the ladder factor is installed in the
commutant rather than derived from a Lorentzian continuum theory.

The positive invariant also defines a genuine finite geometry dynamics.
For a supplied finite language \(\Omega\), positive base weights
\(p_\omega\), and positive interacting operators \(A_\omega\), set
\begin{equation}
 \pi_\omega=
 \frac{p_\omega e^{-g\I_\Delta(A_\omega)}}
 {\sum_{\nu\in\Omega}p_\nu e^{-g\I_\Delta(A_\nu)}} .
 \label{eq:selective-geometry-law}
\end{equation}
On any connected symmetric local move graph, the Metropolis transition
\[
 P_{\omega\nu}=q_{\omega\nu}
 \min\{1,\pi_\nu/\pi_\omega\}
\quad(\nu\neq\omega)
\]
has unique stationary law \(\pi\), obeys exact detailed balance, and admits
the positive-Hilbert isometry
\begin{equation}
 V|\omega\rangle
 =\sum_\nu\sqrt{P_{\omega\nu}}\,
 |\nu\rangle\otimes|e_{\omega\nu}\rangle .
 \label{eq:selective-stinespring}
\end{equation}
Extending \(V\) to a unitary and tracing a fresh environment gives a CPTP
geometry channel.  Hence the finite weights are dynamically sampled rather
than only written as a formal partition sum.

The archived calculation exhausts all 728 connected labelled simple graphs
on five vertices, using
\(A_\omega=0.8\mathbf1+L_\omega\) and \(\Delta=0.72\).
The maximum heat-factorization, detailed-balance, stationary-law, and
Stinespring-normalization residuals are respectively
\(7.00\times10^{-15}\), \(2.71\times10^{-20}\),
\(4.34\times10^{-19}\), and \(2.22\times10^{-16}\).
The exact geometry-channel gap is \(0.1474055064\).  An independent
\(2.0\times10^5\)-sample retained trajectory has total-variation distance
\(0.0248875\) and mean-edge error \(7.24\times10^{-4}\); random Ward
perturbations never exceed one quarter of the analytic distance bound.
These computations test all supplied geometries rather than a selected
saddle.

The finite closure removes the old ``partial interacting Ward protection''
and ``positive graph-weight dynamics'' gaps only at fixed regulator and for
a supplied language.  It does not select \(\Omega\), its topology, or its
move set; it does not give deletion-consistent measures as the regulator
grows; and it does not establish continuum spin--statistics, anomaly
freedom, nonlinear BV closure, or background-independent graph
backreaction.

\subsection{Projective response-weighted causal growth}
\label{sec:projective-response-growth}

The next companion construction
\cite{XuProjectiveCausalGrowth2026} removes one of those qualifications:
the positive response can generate compatible laws at every finite event
number and hence a genuine infinite-history probability measure.  The result
is exact, but it is deliberately stated on a naturally labelled causal-birth
language; quotienting by labellings and taking a Lorentzian continuum limit
remain separate tasks.

Let \(\Omega_n^{(r)}\) be the naturally labelled \(n\)-event partial orders
obtained by adjoining one maximal event at a time.  A birth selects a
down-set \(S\subseteq C\) whose set of maximal elements has size
\[
 p_C(S):=|\max_C S|\leq r
\]
and writes \(C\oplus S\) for the child.  Deleting the largest label gives a
map \(\rho_n:\Omega_{n+1}^{(r)}\to\Omega_n^{(r)}\), and every labelled child
has exactly one parent under this deletion.  On the undirected Hasse graph
of \(C\), define
\begin{equation}
 A_C=m^2\mathbf1+\kappa L_C,\qquad
 \I_\Delta(C)=\Tr f_\Delta(A_C)>0 ,
 \label{eq:pcg-response}
\end{equation}
where
\begin{align}
 f_\Delta(y)
 &=y\log y-2(y+\Delta)\log(y+\Delta)\nonumber\\
 &\quad +(y+2\Delta)\log(y+2\Delta)
 \label{eq:pcg-f}
\end{align}
and, equivalently,
\begin{equation}
 \I_\Delta(C)=
 \int_0^\Delta\!\diff u\int_0^\Delta\!\diff v\,
 \Tr(A_C+u+v)^{-1}.
 \label{eq:pcg-positive}
\end{equation}
This is the \(k=2\) selective-ladder response of
Eq.~\eqref{eq:selective-invariant}, now evaluated on each causal history.

The response of one birth is available from the parent without recomputing
the child spectrum.  Put \(B_C=A_C\oplus m^2\) and
\[
 U_S=\sqrt\kappa\,[e_a-e_{n+1}]_{a\in\max_C S}.
\]
Because a maximal birth changes only its incident Hasse links,
\(A_{C\oplus S}=B_C+U_SU_S^\dagger\).  Woodbury's identity then gives
\begin{align}
 \Delta\I_C(S)
 &:=\I_\Delta(C\oplus S)-\I_\Delta(C)\nonumber\\
 &=f_\Delta(m^2)
 -\int_0^\Delta\!\diff u\int_0^\Delta\!\diff v\,
 \mathcal W_{u+v}(C,S),
 \label{eq:pcg-increment}\\
 \mathcal W_t(C,S)
 &=\Tr\!\left[
 R_tU_S(\mathbf1+U_S^\dagger R_tU_S)^{-1}
 U_S^\dagger R_t\right],
 \qquad R_t=(B_C+t\mathbf1)^{-1}.
 \label{eq:pcg-woodbury}
\end{align}
Thus a \(p\)-parent birth requires only a \(p\times p\) inverse.  Positivity
also yields the nonperturbative parent-number bound
\begin{equation}
 f_\Delta(m^2)-\frac{2\kappa p_C(S)\Delta^2}{m^4}
 \leq\Delta\I_C(S)\leq f_\Delta(m^2).
 \label{eq:pcg-increment-bound}
\end{equation}

For real \(g,\alpha,\beta\), conditionally normalize the allowed births:
\begin{equation}
 G_{n,g}(C,S)=
 \frac{
 \exp[-g\Delta\I_C(S)-\alpha p_C(S)+\beta|S|/n]}
 {\displaystyle\sum_{S'\in\mathcal D_r(C)}
 \exp[-g\Delta\I_C(S')-\alpha p_C(S')+\beta|S'|/n]} .
 \label{eq:pcg-kernel}
\end{equation}
Starting from the one-event order, set
\begin{equation}
 \mu_{n+1,g}(C\oplus S)=\mu_{n,g}(C)G_{n,g}(C,S).
 \label{eq:pcg-recursion}
\end{equation}

\begin{theorem}[Projective response law and positive dilation]
\label{thm:pcg-projective-law}
For every finite \(n\), the laws in
Eq.~\eqref{eq:pcg-recursion} are normalized and satisfy
\begin{equation}
 (\rho_n)_*\mu_{n+1,g}=\mu_{n,g}.
 \label{eq:pcg-projectivity}
\end{equation}
They therefore define a unique probability measure on infinite
past-finite naturally labelled causal histories.  Moreover,
\begin{equation}
 V_{n,g}|C\rangle=
 \sum_{S\in\mathcal D_r(C)}
 \sqrt{G_{n,g}(C,S)}\,
 |C\oplus S\rangle\otimes|e_{C,S}\rangle
 \label{eq:pcg-isometry}
\end{equation}
is an isometry.  Extending it to a unitary and tracing a fresh environment
implements each birth by a CPTP channel.
\end{theorem}

\begin{proof}
Every child has the unique labelled parent obtained by deleting \(n+1\).
Summing Eq.~\eqref{eq:pcg-recursion} over that parent's children and using
\(\sum_SG_{n,g}(C,S)=1\) proves
Eq.~\eqref{eq:pcg-projectivity}.  Kolmogorov extension on the finite
discrete cylinder spaces gives the infinite law.  Orthogonality of the
environment labels and conditional normalization imply
\(V_{n,g}^\dagger V_{n,g}=\mathbf1\).
\end{proof}

The same conditional structure supplies exact observables along a path.
For an \(N\)-event history, let
\[
 X_n=\Delta\I_{C_n}(S_n),\qquad
 \bar X_n=\mathbb E_g[X_n\mid\mathcal F_n].
\]
Its response score is the martingale
\begin{equation}
 \mathcal S_N(g)=\partial_g\log\mathbb P_g(C_1,\ldots,C_N)
 =-\sum_{n=1}^{N-1}(X_n-\bar X_n),
 \label{eq:pcg-score}
\end{equation}
and therefore
\begin{align}
 \mathbb E_g\mathcal S_N&=0,\nonumber\\
 \mathcal F_N(g):=\mathbb E_g\mathcal S_N^2
 &=\sum_{n=1}^{N-1}
 \mathbb E_g\operatorname{Var}_g(X_n\mid\mathcal F_n),\label{eq:pcg-fisher}\\
 \partial_g\mathbb E_g O&=\mathbb E_g[O\mathcal S_N]
 \label{eq:pcg-linear-response}
\end{align}
for every bounded path observable \(O\).  Path relative entropy likewise
obeys the exact chain rule
\begin{equation}
 D(\mathbb P_g^{(N)}\Vert\mathbb P_{g'}^{(N)})
 =\sum_{n=1}^{N-1}
 \mathbb E_gD(G_{n,g}(C_n,\cdot)\Vert
 G_{n,g'}(C_n,\cdot)).
 \label{eq:pcg-kl-chain}
\end{equation}
These identities distinguish a response-generated path law from a
postselected terminal partition sum.

The archived \(r=2\) calculation enumerates every allowed state through
seven events:
\[
 1,\ 2,\ 7,\ 39,\ 322,\ 3730,\ 58113.
\]
The largest normalization, deletion-projectivity, and Stinespring residuals
are \(2.22\times10^{-16}\), \(5.55\times10^{-17}\), and
\(4.44\times10^{-16}\).  The Fisher-chain residual is zero at stored
precision, the relative-entropy-chain residual is
\(2.58\times10^{-17}\), and independent differentiation verifies
Eq.~\eqref{eq:pcg-linear-response} to \(1.84\times10^{-12}\).  Forty-seven
independent birth checks bound the Woodbury and integrated-response errors
by \(2.89\times10^{-16}\) and \(8.89\times10^{-15}\).  A
\(2.5\times10^5\)-history sample reproduces a predeclared 131-bin
trajectory coarse graining with total-variation distance \(0.00677\);
the complete \(58113\)-state histogram is retained rather than hidden by
coarse graining.  As a failure control, independently normalized terminal
Gibbs weights have deletion defect as large as \(0.08635\).

This closes A11PG only at model level.  Natural labels, the precursor
language, \(m^2,\kappa,\Delta,g,\alpha,\beta\), the Hasse response operator,
and the external birth index are inputs.  No equality of different linear
extensions, Rideout--Sorkin discrete general covariance, Bell causality,
coherent quantum measure, dimensional selection, tight Lorentzian
metric--measure continuum limit, nonlinear BV master equation, continuum
matter, graph backreaction, spin--statistics theorem, or anomaly
cancellation is derived.  The result is an exact positive projective
history interface, not a completed nonperturbative gravity measure.

\subsection{Covariant response growth and an extendible quantum measure}
\label{sec:covariant-response-growth}

The preceding conditional kernel solves deletion consistency but not the
natural-label gauge problem.  A second companion construction
\cite{XuCovariantQuantumResponseGrowth2026} identifies a response-generated
subfamily for which the stronger CSG constraints are exact.  The
construction also isolates a coherent subfamily whose cylinder amplitudes
extend countably additively.  These two branches should not be conflated:
the positive branch has a local fresh-environment CPTP dilation, whereas the
coherent branch has a strongly positive rank-one decoherence functional but
no derived local unitary transfer operator.

For each \(k\geq0\), let \(S_k=K_{1,k}\) be the \(k\)-contact star, with
\(S_0=K_1\), and use the same positive logarithmic response as in
Eq.~\eqref{eq:pcg-response}.  Define
\begin{align}
 \chi_k&=\I_\Delta(S_k)-\I_\Delta(S_0),&
 t_k(g)&=e^{-g\chi_k}.
 \label{eq:cqg-contact-couplings}
\end{align}
The star spectrum gives
\begin{equation}
 \chi_0=0,\qquad
 \chi_k=(k-1)f_\Delta(m^2+\kappa)
 +f_\Delta[m^2+\kappa(k+1)]>0
 \label{eq:cqg-star-response}
\end{equation}
for \(k\geq1\).  Thus \(t_k(g)>0\) for every \(k\).  If a precursor
down-set \(S\subseteq C\) has \(v=|S|\) elements and \(m=|\max_C S|\)
maxima, put
\begin{align}
 \lambda_g(v,m)
 &=\sum_{k=m}^{v}\binom{v-m}{k-m}t_k(g),\label{eq:cqg-lambda}\\
 q_{n,g}(C,S)
 &=\frac{\lambda_g(|S|,|\max_C S|)}{\lambda_g(n,0)}.
 \label{eq:cqg-transition}
\end{align}
This transition has a microscopic seed sampler.  Choose
\(K\subseteq C\) with weight
\begin{equation}
 \mathbb P_g(K\mid C)=\frac{t_{|K|}(g)}{\lambda_g(n,0)}
 \label{eq:cqg-seed}
\end{equation}
and place the new maximal event above \(S=\downarrow K\).  The seeds with
down-closure \(S\) are exactly
\(\max S\subseteq K\subseteq S\).  Grouping them by cardinality gives the
binomial coefficients in Eq.~\eqref{eq:cqg-lambda} and proves the Markov
identity
\begin{equation}
 \sum_{S\in\mathcal D(C)}
 \lambda_g(|S|,|\max S|)=\lambda_g(n,0).
 \label{eq:cqg-markov}
\end{equation}

\begin{theorem}[Covariant positive response growth]
\label{thm:cqg-positive-growth}
The transition law in Eq.~\eqref{eq:cqg-transition} is normalized and
deletion-projective.  Its path probability has the intrinsic form
\begin{equation}
 p_g(C)=
 \frac{\displaystyle\prod_{x\in C}
 \lambda_g(|J^-(x)|,|\max J^-(x)|)}
 {\displaystyle\prod_{j=0}^{|C|-1}\lambda_g(j,0)},
 \label{eq:cqg-intrinsic-path}
\end{equation}
and is therefore independent of the natural labelling.  For any two
precursors \(S_1,S_2\) and \(B=S_1\cup S_2\),
\begin{equation}
 q_C(S_1)q_B(S_2)=q_C(S_2)q_B(S_1).
 \label{eq:cqg-bell}
\end{equation}
Thus the law satisfies discrete general covariance and the zero-safe
cross-multiplied Bell condition.  Moreover,
\begin{equation}
 V_{n,g}|C\rangle=
 \sum_{K\subseteq C}
 \sqrt{\frac{t_{|K|}(g)}{\lambda_g(n,0)}}\,
 |C\oplus\downarrow K\rangle|C,K\rangle_E
 \label{eq:cqg-isometry}
\end{equation}
is an isometry and gives a fresh-environment CPTP birth channel.
\end{theorem}

\begin{proof}
Equation~\eqref{eq:cqg-markov} proves normalization and the unique labelled
parent proves deletion projectivity.  Multiplication along a path gives
Eq.~\eqref{eq:cqg-intrinsic-path}; its numerator is a product of intrinsic
past invariants and its denominator depends only on cardinality.  Removing
the spectators \(C\setminus B\) leaves the two precursor numerators
unchanged, while the row normalizer cancels in
Eq.~\eqref{eq:cqg-bell}.  Finally,
\(\sum_{K\subseteq C}t_{|K|}=\lambda_g(n,0)\) proves
\(V_{n,g}^\dagger V_{n,g}=\mathbf1\).
\end{proof}

The response remains operational along the covariant paths.  Define
\begin{equation}
 R_g(v,m)=-\partial_g\log\lambda_g(v,m)
 =\frac{\sum_{k=m}^{v}\binom{v-m}{k-m}t_k(g)\chi_k}
 {\lambda_g(v,m)}.
 \label{eq:cqg-response-mean}
\end{equation}
The score increment of a birth is
\[
 Y_n=R_g(n,0)-R_g(v_n,m_n).
\]
Differentiating Eq.~\eqref{eq:cqg-markov} gives
\(\mathbb E_g[Y_n\mid\mathcal F_n]=0\), so the Fisher, linear-response, and
path-relative-entropy chain rules in
Eqs.~\eqref{eq:pcg-fisher}--\eqref{eq:pcg-kl-chain} hold with \(Y_n\) in
place of the global action increment.

The coherent branch retains one response-generated contact amplitude:
\begin{equation}
 \widetilde t_0=1,\qquad
 \widetilde t_r=z_r=e^{-g\chi_r/2+i\theta\chi_r},\qquad
 \widetilde t_k=0\quad(k\neq0,r).
 \label{eq:cqg-complex-coupling}
\end{equation}
The complex version of Eqs.~\eqref{eq:cqg-lambda} and
\eqref{eq:cqg-transition} has row sum one, the same intrinsic product, and
the cross-multiplied Bell identity even when some transitions vanish.
Let
\[
 \zeta_n(C)=\sum_{S\in\mathcal D(C)}
 |\widetilde q_n(C,S)|-1,\qquad
 \zeta_n^{\max}=\max_{|C|=n}\zeta_n(C).
\]
For \(s=|z_r|\) and \(\phi=\arg z_r\), the maximum occurs at the
\(n\)-antichain \cite{SuryaZalel2020} and equals
\begin{equation}
 \zeta_n^{\max}
 =\frac{1+\binom nr s}
 {\sqrt{1+2s\binom nr\cos\phi+s^2\binom nr^2}}-1.
 \label{eq:cqg-zeta}
\end{equation}

\begin{theorem}[Extendible coherent response history]
\label{thm:cqg-complex-extension}
If \(r\geq2\), \(g>0\), and
\(\phi=\theta\chi_r\in(0,\pi/2]\), the complex cylinder measure generated by
Eq.~\eqref{eq:cqg-complex-coupling} has bounded variation and extends
uniquely to a countably additive complex measure \(\mu\) on the cylinder
sigma algebra.  The functional
\begin{equation}
 \mathfrak D(A,B)=\mu(A)\overline{\mu(B)}
 \label{eq:cqg-decoherence}
\end{equation}
is normalized, countably additive, and strongly positive; its diagonal
\(\mathfrak q(A)=|\mu(A)|^2\) obeys the grade-2 quantum-measure sum rule.
For \(r=1\) and nonreal \(z_1\), bounded variation fails.
\end{theorem}

\begin{proof}
Put \(x=s\binom nr\).  Rationalizing Eq.~\eqref{eq:cqg-zeta} gives
\[
 \zeta_n^{\max}
 =\frac{2x(1-\cos\phi)}
 {\sqrt{1+2x\cos\phi+x^2}
 [1+x+\sqrt{1+2x\cos\phi+x^2}]}
 \leq\frac{1-\cos\phi}{s\binom nr}.
 \]
The upper bound is summable for \(r\geq2\), so the finite cylinder
variations are bounded by
\(\prod_n(1+\zeta_n^{\max})<\infty\).  The complex-measure extension
criterion of Ref.~\cite{SuryaZalel2020} gives \(\mu\).  Every finite matrix
\([\mathfrak D(A_i,A_j)]\) is the outer product
\([\mu(A_i)][\mu(A_j)]^\dagger\), hence is positive semidefinite.  Ordinary
additivity of \(\mu\) yields the grade-2 sum rule
\cite{Sorkin1994QuantumMeasure}.  For nonreal \(r=1\),
\(\zeta_n^{\max}\sim(1-\cos\phi)/(sn)\), and the variation diverges.
\end{proof}

The coherent measure is not merely a complex rephrasing of a probability
law.  For \(r=2\), the two allowed three-event unlabelled endpoint classes
are the antichain \(A_3\) and the two-parent order \(V_3\), with amplitudes
\[
 \mu_3(A_3)=\frac1{1+z_2},\qquad
 \mu_3(V_3)=\frac{z_2}{1+z_2}.
 \]
Their finite-volume quotient interference is
\begin{equation}
 I_3=
 \mathfrak q(A_3\cup V_3)-\mathfrak q(A_3)-\mathfrak q(V_3)
 =\frac{2\Re z_2}{|1+z_2|^2}.
 \label{eq:cqg-interference}
\end{equation}
For the archived benchmark, \(I_3=0.3312668\).

Complete enumeration gives
\[
 1,\ 2,\ 7,\ 40,\ 357,\ 4824
\]
naturally labelled histories and
\[
 1,\ 2,\ 5,\ 16,\ 63,\ 318
\]
unlabelled classes through six events.  Markov, seed-grouping, intrinsic
path, DGC, Bell, Fisher, and relative-entropy residuals are all below
\(1.8\times10^{-15}\).  The \(r=2,3\) accumulated row excess stabilizes
with the predicted \(N^{-r}\) tail through \(N=200000\), while the
\(r=1\) control grows logarithmically.

There is also a decisive falsification.  If the full Hasse response is used
as a conditionally normalized global potential,
\[
 q^{\rm glob}_{n,g}(C,S)\propto
 e^{-g[\I_\Delta(C\oplus S)-\I_\Delta(C)]},
 \]
two natural labellings of one five-event causal set differ in log path
weight by \(1.624733\), and removing one spectator changes a competing log
transition ratio by \(0.136424\).  Positivity and projectivity therefore do
not imply DGC or Bell causality; the response must enter through a
covariant contact sequence (or another structure satisfying the same
constraints).

This closes A11CG at model level.  It removes the natural-label and Bell
defects for the response-generated CSG subfamily and supplies a
countably-additive strongly positive coherent history measure.  It does not
select the maximal-birth language, contact stars, response multiplet,
parameters, arity, or phase.  The coherent decoherence functional is rank
one and has no derived local unitary complex transfer.  No Lorentzian
continuum tightness, dimension selection, nonlinear gravitational
constraint algebra, continuum matter, anomaly freedom, or graph
backreaction is proved.

\subsection{A local-unitary non-Abelian history measure on a fixed causal
strip}
\label{sec:nonabelian-history-measure}

The rank-one boundary above can be removed without sacrificing countable
extension, but presently only after fixing a smaller causal language
\cite{XuNonAbelianQuantumHistory2026}.  Let
\[
 \Omega_2=\{00,01,10,11\}^{\mathbb N}
\]
be the record space of a ranked two-port strip.  A letter
\(x_n=(a_n,b_n)\) tells the two ports at rank \(n\) which of the two
ports at rank \(n-1\) is their parent.  Every generated edge raises rank,
so acyclicity is kinematic.  This is not an unlabeled causal-set growth law.

On a two-dimensional bond space, for any Pauli involution \(H\), define
\[
 F_H^{0,1}(\epsilon)=\frac{\mathbf 1\pm e^{i\epsilon H}}{2},
 \qquad
 A_n^{ab}=F_{H_a}^{b}(\epsilon_n)F_X^a(\epsilon_n),
 \qquad (H_0,H_1)=(Z,Y).
\]
The angle is supplied by the same positive contact response,
\[
 \epsilon_n=\epsilon_\star e^{-g\chi_n}/n^p,
 \qquad 0<\epsilon_\star<\pi.
\]
The two binary decisions have a Hadamard--controlled-unitary--Hadamard
dilation with fresh ancillas.  Coherently conditioning the second axis on
the first record gives all four blocks.

\begin{theorem}[Local-unitary non-Abelian history extension]
\label{thm:nonabelian-history-extension}
For every cell,
\begin{equation}
 \sum_{a,b}A_n^{ab}=\mathbf 1,
 \qquad
 \sum_{a,b}A_n^{ab\dagger}A_n^{ab}=\mathbf 1.
 \label{eq:nonabelian-history-two-sums}
\end{equation}
The transition algebra is noncommutative:
\begin{equation}
 [A^{00},A^{10}]=\frac{i}{2}\epsilon^2Y+O(\epsilon^3).
 \label{eq:nonabelian-history-commutator}
\end{equation}
For a cylinder word \(w=x_1\cdots x_N\), set
\[
 M([w])=A_N^{x_N}\cdots A_1^{x_1}.
\]
If \(\sum_n\epsilon_n<\infty\), this operator cylinder content extends
uniquely to a countably additive regular
\(M_2(\mathbb C)\)-valued measure on the Borel sigma algebra of
\(\Omega_2\).  For every density matrix \(\rho\),
\begin{equation}
 D(E,F)=\Tr[\rho\,M(E)^\dagger M(F)]
 \label{eq:nonabelian-history-decoherence}
\end{equation}
is normalized and strongly positive, and \(D(E,E)\) is a grade-2 quantum
measure.  Two independent strips obey exact spectator factorization.
\end{theorem}

\begin{proof}
For either binary decision,
\[
 F_H^0+F_H^1=\mathbf 1,\qquad
 F_H^{0\dagger}F_H^0+F_H^{1\dagger}F_H^1=\mathbf 1.
\]
Summing first over \(b\) and then \(a\) gives
Eq.~\eqref{eq:nonabelian-history-two-sums}; the circuit above proves local
unitary realizability.  The Pauli expansion gives
Eq.~\eqref{eq:nonabelian-history-commutator}.  Cylinder refinement is exact
because the first identity in
Eq.~\eqref{eq:nonabelian-history-two-sums} is an operator Markov sum.
Moreover,
\begin{align}
 \sum_{|w|=N}\|A_w\|
 &\le
 \prod_{n=1}^{N}
 \left(\cos\frac{\epsilon_n}{2}
      +\sin\frac{\epsilon_n}{2}\right)^2\nonumber\\
 &\le \exp\!\left(\sum_n\epsilon_n\right).
 \label{eq:nonabelian-history-variation}
\end{align}
Thus every matrix entry defines a uniformly bounded functional on the
locally constant functions of the compact product space.  Density in
\(C(\Omega_2)\) followed by Riesz--Markov gives the unique countably
additive matrix measure.  Equation
\eqref{eq:nonabelian-history-decoherence} is a Hilbert--Schmidt Gram
functional, which proves strong positivity and the grade-2 identity.
Tensor-product instruments and product bond states give spectator
factorization after summing the complete spectator outcome partition.
\end{proof}

The response sequence is summable for every \(p\ge0\) when \(g>0\),
because \(\chi_n\ge(n-1)f_\Delta(m^2+\kappa)>0\); at \(g=0\), \(p>1\)
is sufficient and necessary for this power-law tail.  At the archived
benchmark \(\epsilon=0.7\),
\[
 2\Re D([00],[10])=-\frac14\sin^3\epsilon=-0.06684023094.
\]
The history Gram rank is three after one cell and reaches the full
Hilbert--Schmidt bond rank four after two; the transition commutator is
\(0.1949275542\).  The largest Markov, Kraus, dilation, and spectator
residual is \(6.3\times10^{-16}\).  A commuting-axis control gives zero
commutator, while the \(g=0,p=1\) variation envelope grows to
\(1.7604\times10^4\) by \(N=200000\).  The latter is failure of the
sufficient extension certificate, not a proof of nonextension.

This proves A11NH at model level.  It removes the scalar/rank-one and
local-realizability restrictions for one fixed causal strip.  It does not
combine that non-Abelian transfer with the unlabeled DGC and Bell-causal
arbitrary-down-set growth of A11CG.  The strip, rank foliation, port
language, bond dimension, response sequence, and summable tail are inputs.
This restriction is substantive: natural operator orderings of quantum
Bell causality on unrestricted sequential growth can force the transition
algebra to commute, and a general causal-past-ordered representation remains
unknown \cite{SrivastavaSurya2026}.
No fluctuating generic adjacency, Lorentzian continuum tightness, nonlinear
constraint algebra, continuum matter, or graph backreaction is proved.

\subsection{A covariant non-Abelian operator measure on arbitrary causal
histories}
\label{sec:covariant-nonabelian-holonomy}

The fixed-strip restriction can be removed without giving up countable
extension if the noncommuting operator is treated as an event-local density
over the covariant positive growth law, rather than as a square transition
ratio \cite{XuCovariantNonAbelianHolonomy2026}.  Use the response-generated
CSG probability \(\mathbb P_g\) of
Theorem~\ref{thm:cqg-positive-growth}.  If an event \(x\) is born above a
precursor whose volume and number of maximal elements are
\[
 v_x=|J^-(x)|,\qquad m_x=|\max J^-(x)|,
\]
attach a fresh qubit and apply
\begin{equation}
 u_x=
 \exp\!\left[
 i\epsilon_\star e^{-h\chi_{v_x}}H_{m_x\bmod3}
 \right],
 \qquad (H_0,H_1,H_2)=(X,Y,Z).
 \label{eq:cnh-gate}
\end{equation}
Every factor acts only on the new event qubit.  The aggregate
\((v_x,m_x)\) is intrinsic to the precursor and is unchanged when spectators
outside a competing precursor union are deleted.

Let \(a=f_\Delta(m^2+\kappa)\),
\[
 r=e^{-ga},\qquad s=e^{-ha},\qquad
 \rho=\frac{1+rs}{1+r}<1.
\]
Equation~\eqref{eq:cqg-star-response} gives
\(\chi_k=(k-1)a+b_k\), with \(0<b_k\le a\).  Hence
\begin{equation}
 r^k\le t_k\le e^{ga}r^k,\qquad
 \epsilon_\star e^{-h\chi_k}
 \le \epsilon_\star e^{ha}s^k.
 \label{eq:cnh-geometric-comparison}
\end{equation}
At stage \(n\), the CSG seed cardinality has conditional distribution
\[
 \mathbb P_g(|K_n|=k\mid\mathcal F_n)
 =\binom nk t_k/\lambda_g(n,0).
\]
Since the precursor volume \(V_n=|\downarrow K_n|\) is at least
\(|K_n|\), Eq.~\eqref{eq:cnh-geometric-comparison} gives the
state-uniform tail
\begin{equation}
 \mathbb E_g(s^{V_n}\mid\mathcal F_n)
 \le e^{ga}\rho^n,\qquad
 \mathbb E_g\theta_{V_n}
 \le\epsilon_\star e^{(g+h)a}\rho^n.
 \label{eq:cnh-tail}
\end{equation}
Define
\[
 B_{\rm hol}=
 \frac{\epsilon_\star e^{(g+h)a}}{1-\rho}.
\]

\begin{theorem}[Covariant non-Abelian causal-holonomy measure]
\label{thm:covariant-nonabelian-holonomy}
Suppose \(B_{\rm hol}<1\).  On the incomplete event tensor product
\[
 \mathcal H_{\rm ev}
 =\bigotimes_{x\in\mathbb N}(\mathbb C^2,|0\rangle_x),
\]
the finite products \(U_N(\omega)=\bigotimes_{x\le N}u_x(\omega)\)
converge in operator norm for \(\mathbb P_g\)-almost every causal history to
a unitary \(U(\omega)\).  The Bochner integral
\begin{equation}
 \mathsf M(E)=\int_EU(\omega)\,\mathbb P_g(d\omega)
 \label{eq:cnh-operator-measure}
\end{equation}
is a countably additive \(B(\mathcal H_{\rm ev})\)-valued measure of total
variation at most one.  For every permutation-invariant density operator
\(\varrho\),
\begin{equation}
 D(E,F)=
 \frac{\Tr[\varrho\,\mathsf M(E)^\dagger\mathsf M(F)]}
 {\Tr[\varrho\,\mathsf M(\Omega)^\dagger\mathsf M(\Omega)]}
 \label{eq:cnh-decoherence}
\end{equation}
is normalized and strongly positive, and \(D(E,E)\) is grade 2.  Natural
relabelings act equivariantly,
\[
 U(\pi\omega)=\Pi_\pi U(\omega)\Pi_\pi^\dagger,
\]
and the range of \(\mathsf M\) is noncommutative.
\end{theorem}

\begin{proof}
Equation~\eqref{eq:cnh-tail} and Tonelli's theorem give
\[
 \mathbb E_g\sum_n\theta_{V_n}\le B_{\rm hol}<\infty,
\]
so the angle sum is finite almost surely.  Since
\(\|u_x-\mathbf1\|\le\theta_{V_x}\), telescoping makes \(U_N\)
operator-norm Cauchy.  Its finite-cylinder approximants have separable
range, which proves Bochner measurability.  A bounded Bochner density
defines an operator-norm countably additive vector measure and
\(|\mathsf M|(\Omega)\le\int\|U\|d\mathbb P_g=1\).  Furthermore,
\[
 \|\mathsf M(\Omega)-\mathbf1\|
 \le\mathbb E_g\|U-\mathbf1\|
 \le B_{\rm hol}<1,
\]
so the denominator in Eq.~\eqref{eq:cnh-decoherence} is at least
\((1-B_{\rm hol})^2\).  The Gram identity
\[
 \sum_{ij}\bar c_i c_jD(E_i,E_j)
 \propto
 \Tr\!\left[
 \varrho\left(\sum_i c_i\mathsf M(E_i)\right)^\dagger
 \left(\sum_j c_j\mathsf M(E_j)\right)
 \right]\ge0
\]
proves strong positivity, and biadditivity gives the grade-2 sum rule.
Intrinsic past data prove equivariance under natural relabeling.

For noncommutativity, fix a positive-probability \(n\)-event prefix and
compare a next birth with empty precursor to one with a singleton precursor.
Their fresh-factor gates \(u_{0,0}\) and \(u_{1,1}\) obey
\[
 c_{\rm NA}:=\|[u_{0,0},u_{1,1}]\|
 =2|\sin\theta_0\sin\theta_1|>0.
\]
After integrating unrestricted futures, both conditional tail operators
are within \(B_{\rm hol}\rho^{n+1}\) of the identity.  The corresponding
cylinder commutator is bounded below, apart from strictly positive path
probabilities, by
\begin{equation}
 c_{\rm NA}-4B_{\rm hol}\rho^{n+1}.
 \label{eq:cnh-range-bound}
\end{equation}
This is positive for all sufficiently large \(n\), so integration does not
abelianize the range.
\end{proof}

For the archived parameters
\[
 (m^2,\Delta,\kappa,g,h,\epsilon_\star)
 =(0.85,0.65,0.60,0.82,2.20,0.040),
\]
\(\rho=0.833981\) and \(B_{\rm hol}=0.446848\).
The lower bound~\eqref{eq:cnh-range-bound} becomes positive at \(n=36\).
Complete enumeration through six events gives 4824 naturally labelled
histories in 318 unlabeled classes.  The Markov, scalar Bell, and
spectator-gate residuals are below \(6.7\times10^{-16}\), while intrinsic
decorated-causet covariance is exact.  A deliberately stage-labelled gate
has covariance and spectator defects \(2.34\times10^{-2}\) and
\(1.00\times10^{-2}\).  At four events the normalized decoherence matrix
has interference \(1.10\times10^{-2}\), grade-2 residual
\(5.4\times10^{-19}\), and a nonzero operator-measure commutator.

This proves A11CH at model level.  It combines arbitrary-down-set CSG
covariance with a local-unitary non-Abelian, countably additive history
measure.  The exact spectator statement is
\begin{equation}
 \frac{q_C(S_1)}{q_C(S_2)}
 =\frac{q_B(S_1)}{q_B(S_2)},\qquad
 u_C(S_i)=u_B(S_i),
 \quad B=S_1\cup S_2.
 \label{eq:cnh-isometric-bell}
\end{equation}
It is \emph{isometric Bell locality}, not the square-operator TOBC, NTOBC,
or CPOBC cross-product condition.  The operator density is integrated
against a classical CSG law and does not obey an operator Markov sum at each
parent.  Geometry does not backreact on the gate, precursor locality is not
bounded-neighborhood locality, and no Lorentzian continuum, nonlinear
constraint algebra, matter dynamics, anomaly freedom, or graph backreaction
is proved.

\subsection{Central operator-valued growth and the no-coherence boundary}
\label{sec:central-operator-growth}

The preceding constructions leave a tempting shortcut: promote the scalar
CSG couplings to positive operators and interpret their state dependence as
quantum backreaction.  The commutative branch of that proposal can be solved
and classified completely \cite{XuOperatorValuedCausalGrowth2026}.  The
classification is useful precisely because it proves that central
operator-valued couplings do \emph{not} make the geometry marginal
coherent.

Let \(\mathcal Z\) be a commutative von Neumann algebra with separable
predual,
\[
 \mathcal Z\simeq L^\infty(Z,\nu),
 \qquad T_k=M_{t_k}\in\mathcal Z_+,\qquad T_0=\mathbf1 .
\]
For a birth above a down-set \(S\) of an \(n\)-event parent \(C\), put
\[
 v=|S|,\qquad m=|\max S|,
\]
and define
\begin{align}
 \Lambda(v,m)
 &=\sum_{k=m}^{v}\binom{v-m}{k-m}T_k,\nonumber\\
 \Lambda_n&=\Lambda(n,0),\qquad
 Q_C(S)=\Lambda(v,m)\Lambda_n^{-1}.
 \label{eq:ocg-transition}
\end{align}
The normalization is an operator identity rather than a fitted row sum.
Every seed \(K\subseteq C\) has the unique down-closure
\(\downarrow K=S\), and the seeds in that fiber are exactly
\(\max S\subseteq K\subseteq S\).  Hence
\begin{equation}
 \sum_{S\in\mathcal D(C)}
 \Lambda(|S|,|\max S|)
 =\sum_{K\subseteq C}T_{|K|}
 =\Lambda_n .
 \label{eq:ocg-seed-identity}
\end{equation}
Since \(T_0=\mathbf1\), \(\Lambda_n\geq\mathbf1\); therefore
\(Q_C(S)\geq0\) and \(\sum_SQ_C(S)=\mathbf1\).

\begin{theorem}[Central operator-valued growth and classification]
\label{thm:central-operator-growth-classification}
The effects in Eq.~\eqref{eq:ocg-transition} have the following
properties.
\begin{enumerate}
\item \(K_C(S)=Q_C(S)^{1/2}\) is a trace-preserving L\"uders instrument at
each finite parent.
\item For every finite history \(\omega\),
\begin{equation}
 K_\omega=\prod_jK_{C_j}(S_j)
 =M_{\sqrt{P_z(\omega)}} ,
 \label{eq:ocg-path}
\end{equation}
where \(P_z\) is the scalar CSG law with couplings \(t_k(z)\).  Thus
naturally labeled paths with isomorphic endpoints have identical path
operators.
\item For \(B=S_1\cup S_2\), the exact Bell cross-product is
\begin{equation}
 Q_C(S_1)Q_B(S_2)=Q_C(S_2)Q_B(S_1).
 \label{eq:ocg-bell}
\end{equation}
\item The finite effects extend to a normalized strongly countably
additive POVM on the infinite past-finite history event algebra,
\begin{equation}
 Q(E)=M_{P_z(E)}.
 \label{eq:ocg-infinite-povm}
\end{equation}
\item Conversely, every central operator-valued CSG has this direct-integral
form.  For every normal state \(\rho\),
\begin{equation}
 p_\rho(E)=\Tr[\rho Q(E)]
 =\int_ZP_z(E)\,\diff\mu_\rho(z),
 \label{eq:ocg-mixture}
\end{equation}
where \(\mu_\rho\) is fixed by the restriction of \(\rho\) to
\(\mathcal Z\).
\end{enumerate}
Consequently two states with the same central restriction give identical
geometry statistics on every history event.
\end{theorem}

\begin{proof}
Equation~\eqref{eq:ocg-seed-identity} proves Kraus completeness.  In the
multiplication representation,
\[
 Q_C(S)=M_{q_z(C,S)},\qquad
 q_z(C,S)=
 \frac{\sum_{k=m}^{v}\binom{v-m}{k-m}t_k(z)}
 {\sum_{k=0}^{n}\binom nk t_k(z)} .
\]
The scalar path product depends only on the intrinsic past
\((|J^-(x)|,|\max J^-(x)|)\) of each endpoint event, proving
Eq.~\eqref{eq:ocg-path} and DGC.  The two precursor numerators are unchanged
when spectators outside \(B\) are removed, while the parent normalizer is
common; commutativity then gives Eq.~\eqref{eq:ocg-bell}.  The scalar
kernels define probability measures \(P_z\) on infinite histories.
Measurability first holds for cylinders and extends by the monotone-class
theorem.  Scalar countable additivity and dominated convergence give strong
countable additivity of Eq.~\eqref{eq:ocg-infinite-povm}.  Finally, the
representation theorem for commutative von Neumann algebras makes every
central coupling a multiplication function, and every normal functional on
\(\mathcal Z\) is integration against \(\mu_\rho\), which proves
Eq.~\eqref{eq:ocg-mixture}.
\end{proof}

For two central sectors,
\[
 Q(E)=
 \begin{pmatrix}P_0(E)&0\\0&P_1(E)\end{pmatrix},\qquad
 \rho=
 \begin{pmatrix}w&c\\\bar c&1-w\end{pmatrix},
\]
so
\begin{equation}
 p_\rho(E)=wP_0(E)+(1-w)P_1(E)
 \label{eq:ocg-coherence-blind}
\end{equation}
is independent of \(c\).  Even when an external reference makes the
off-diagonal block operational in a larger algebra, the central geometry
POVM cannot detect it.  The geometry can be state controlled or classically
correlated with a hidden sector, but it cannot interfere between sectors.

Positivity and one-step normalization alone do not evade this conclusion.
For noncommuting \(T_k>0\), the symmetrically normalized effects
\begin{equation}
 E_C(S)=
 \Lambda_n^{-1/2}\Lambda(v,m)\Lambda_n^{-1/2}
 \label{eq:ocg-noncentral-effect}
\end{equation}
remain positive and sum to the identity, yet their square roots need not
compose covariantly.  An explicit two-dimensional control is
\begin{equation}
 T_k=e^{-\bar g\chi_k}
 \left[
 \mathbf1+\eta\cos(0.79k)X+\eta\sin(0.79k)Z
 \right],
 \qquad 0<\eta<1 ,
 \label{eq:ocg-noncentral-coupling}
\end{equation}
with the positive contact-star response \(\chi_k\) of
Eq.~\eqref{eq:cqg-star-response}.  For the archived parameters, the two
natural labelings
\[
 (0,1,0),\qquad (0,0,1)
\]
of a two-event chain plus one isolated event have sequential probabilities
\[
 0.14917364759401094,\qquad
 0.14919932629598825.
\]
Their DGC spread is \(2.5678702\times10^{-5}\), although the maximum
parentwise POVM residual is \(1.11\times10^{-15}\).  This is a failure
control, not a proposed noncentral growth law.

Non-Abelian records can nevertheless coexist with the central geometry
marginal.  For \(U_H=e^{i\epsilon H}\), set
\[
 F_H^0=\frac{\mathbf1+U_H}{2},\qquad
 F_H^1=\frac{\mathbf1-U_H}{2},\qquad
 A^{ab}=F_{H_a}^{\,b}F_X^{\,a},
 \quad(H_0,H_1)=(Z,Y).
\]
The two binary identities imply
\[
 \sum_{ab}A^{ab}=\mathbf1,\qquad
 \sum_{ab}A^{ab\dagger}A^{ab}=\mathbf1,
\]
while at \(\epsilon=0.22\),
\[
 \|[A^{00},A^{10}]\|=2.0800387\times10^{-2}.
\]
The combined fresh-event instrument
\begin{equation}
 L_C(S,a,b)=Q_C(S)^{1/2}\otimes A^{ab}
 \label{eq:ocg-combined-instrument}
\end{equation}
is trace preserving, but summing over \(a,b\) returns \(Q_C(S)\); the
record noncommutativity does not backreact on the geometry law.

Complete enumeration through six events covers 4824 naturally labeled
histories in 318 unlabeled classes.  The largest central POVM, DGC, and
Bell residuals are \(6.66\times10^{-16}\),
\(2.78\times10^{-17}\), and \(2.78\times10^{-17}\); direct-mixture and
coherence-visibility residuals vanish at stored precision.  The analytic
classification, rather than these finite residuals, is the infinite-history
result.

This proves A11OC at model level and sharpens the remaining cut.  The
infinite object established here is a POVM, not a countably additive
noncommutative instrument after unobserved future records are discarded.
No noncentral solution of causal-past-ordered Bell causality, coherent
geometry interference, geometry backreaction, tight Lorentzian continuum,
constraint algebra, anomaly theorem, or nonlinear Einstein limit is
obtained.  A viable quantum-growth replacement must therefore leave the
central algebra while preserving normalization, DGC, and operator locality,
or must use a different strongly positive history construction whose
relative phases survive both quotienting and infinite extension.

\subsection{Rigidity of the remaining square-operator Bell branch}
\label{sec:bell-causal-rigidity}

The preceding classification does not decide whether noncentral transition
operators can obey an operator Markov sum, operator discrete general
covariance (DGC), and causal-past-ordered Bell causality (CPOBC)
simultaneously.  CPOBC is the only nonsingular ordering in the current
square-operator formulation that is not already known to force immediate
commutativity \cite{SrivastavaSurya2026}.  The companion analysis
\cite{XuBellCausalGrowthRigidity2026} removes three natural neighborhoods
from that search.

Let \(a_n\) be the \(n\)-element antichain and set
\[
 Q_n=A(a_n,\varnothing).
\]
Pairing the gregarious transition with a singleton precursor and deleting
their common spectators gives the exact CPOBC identity
\begin{equation}
 A_n^{(1)}Q_1=(\mathbf1-Q_1)Q_n,
 \qquad
 A_n^{(1)}=(\mathbf1-Q_1)Q_nQ_1^{-1}.
\label{eq:bcgr-one-precursor}
\end{equation}
If the three displayed transition operators are self-adjoint, then
\begin{equation}
 Q_1\!\left(A_n^{(1)}-A_n^{(1)\dagger}\right)Q_1
 =
 [\,Q_1(\mathbf1-Q_1),Q_n\,].
\label{eq:bcgr-hermiticity}
\end{equation}
Thus self-adjointness makes \(Q_n\) commute with
\(Q_1(\mathbf1-Q_1)\).  The
identity alone has a complementary-spectrum loophole, because the map
\(f(\lambda)=\lambda(1-\lambda)\) identifies
\(\lambda\) and \(1-\lambda\).  The full CPOBC system closes that loophole.

\begin{theorem}[Self-adjoint and all-order triangular CPOBC rigidity]
\label{thm:bcgr-rigidity}
For nonsingular square-operator CPOBC:
\begin{enumerate}
\item if every transition is self-adjoint, then the complete
finite-dimensional transition algebra is commutative, with no positivity
or sign assumption and including all complementary-spectrum blocks;
\item through four events, at the scalar transitive-percolation point
\(t=1/4\), the 79 operator residuals in 30 marked-transition classes have
\begin{equation}
 \operatorname{rank}J=27,\qquad
 \operatorname{rank}[\,J\mid W_{12}\ W_{13}\ W_{23}\,]=30,
\label{eq:bcgr-ranks}
\end{equation}
so the three matrices multiplying every leading formal deformation obey
\([H_1,H_2]=[H_1,H_3]=[H_2,H_3]=0\);
\item for every pair of distinct strictly positive scalar CSG coupling
sequences \(\bm t_L,\bm t_R\in(0,\infty)^{\mathbb N}\), every
\(2\times2\) upper-triangular extension
\[
 A_e=
 \begin{pmatrix}
 a_e(\bm t_L)&x_e\\0&a_e(\bm t_R)
 \end{pmatrix},
\]
is globally similar to the diagonal direct sum; if \(t_k\) is the first
unequal coupling, the universal early obstruction has
\begin{equation}
\det\mathcal O_k=
\frac{t_2}{\Lambda_1\Lambda_2}
\prod_{j=2}^{k-1}\frac{t_1t_{j-1}}{\Lambda_j^2}>0
\qquad(k\ge3).
\label{eq:bcgr-allorder-det}
\end{equation}
Independently, the six-event \(9680\times707\) linear system has rank 706
on every distinct positive five-jet and its kernel is only
\[
 \bm x=\xi\,[\,\bm a(\bm t_L)-\bm a(\bm t_R)\,],
\]
the coboundary generated by a global similarity transform.
\item for coincident positive scalar characters, every \(2\times2\)
upper-triangular self-extension is uniquely
\begin{equation}
 A_e=a_e(\bm t)\mathbf1+
 \left(\sum_{\ell\ge1}c_\ell\partial_{t_\ell}a_e(\bm t)\right)E_{12},
\label{eq:bcgr-equal-character}
\end{equation}
where a finite transition contains only finitely many terms.  Its
transition algebra is commutative, and it splits if and only if all
\(c_\ell\) vanish.
\item if a nonsingular \(2\times2\) CPOBC representation is irreducible
and \(R_n=Q_1^{-1}Q_n\), then \(R_2\) is scalar.  A non-scalar \(R_2\)
reduces to an eigenline-swapping normal form, whose two exceptional
components are both eliminated by the following antichain stage.
\item write the surviving scalar ratio as \(R_2=h\mathbf1\).  The
two-element chain requires \(h\ne0,1\).  In the trace-balanced subchart
\(d=a,\ x=-1\) of the first-simple-spectrum chart
\begin{equation}
 Q_1=\begin{pmatrix}a&1\\q&d\end{pmatrix},\quad
 R_3=\operatorname{diag}(1,x),\quad
 R_4=\operatorname{diag}(1,y)
\label{eq:bcgr-delayed-component}
\end{equation}
the physically saturated stage-three equations have exactly
\[
C_0:y=1,\quad
E_1:a=\tfrac12,\ 4hq-3h+2=0,\quad
E_2:q=a-a^2,\ 2ah-2h+1=0.
\]
None has a nonsingular extension through the next antichain stage.  No
classification of the complementary part of the first-simple-spectrum
chart is asserted.
\end{enumerate}
\end{theorem}

\begin{proof}
Put \(F=Q_1(\mathbf1-Q_1)\).  Equation~\eqref{eq:bcgr-hermiticity}
gives \([F,Q_n]=0\).  A second exact CPOBC relation is
\begin{equation}
 [Q_1Q_2^{-1},Q_1^{-1}Q_2]=0.
\label{eq:bcgr-q1q2}
\end{equation}
The following elementary block argument removes the former sign
restriction.  Let \(A,B\) be invertible self-adjoint matrices satisfying
\[
 [A(\mathbf1-A),B]=0,\qquad [AB^{-1},A^{-1}B]=0 .
\]
On a spectral subspace of \(A(\mathbf1-A)\), the spectrum of \(A\) contains
at most the complementary pair \(q,1-q\).  In the nontrivial case write
\[
 A=\begin{pmatrix}q\mathbf1&0\\0&(1-q)\mathbf1\end{pmatrix},
 \qquad
 B=\begin{pmatrix}U&D\\D^\dagger&V\end{pmatrix}.
\]
With \(C=A^{-1}B\), the second commutator is equivalent to
\([C,ACA^{-1}]=0\).  Its upper diagonal block is
\[
 \bigl[q^{-2}-(1-q)^{-2}\bigr]DD^\dagger .
\]
Invertibility excludes \(q=0,1\); for \(q\ne1/2\) the coefficient is
nonzero and hence \(D=0\), while at \(q=1/2\) the two eigenvalues already
coincide.  Thus \([A,B]=0\).  Applying this lemma to \(Q_1,Q_2\) proves
\([Q_1,Q_2]=0\).

Set \(R=Q_2Q_1^{-1}\).  It is self-adjoint and commutes with \(Q_1\).
The three transitions from the two-element chain reduce to
\begin{equation}
 G_2=Q_2,\qquad
 V_2=(\mathbf1-Q_1)R,\qquad
 T_2=\mathbf1-R .
\label{eq:bcgr-chain}
\end{equation}
Nonsingularity of \(T_2\) gives \(1\notin\operatorname{spec}R\), while the
three-generator CPOBC relation gives \([R,Q_n]=0\).

It remains to exclude mixing of complementary \(Q_1\) eigenvalues.  The
two-element-antichain timid transition is
\[
 S=A_2^{(2)}=\mathbf1-(2\mathbf1-Q_1)R,
 \qquad
 A_n^{(2)}=S Q_nQ_2^{-1}.
\]
Resolve the commuting pair \((Q_1,R)\) into joint spectral projections
\(P_{q,r}\).  The relations \([F,Q_n]=[R,Q_n]=0\) permit an off-diagonal
block only between \((q,r)\) and \((1-q,r)\).  Hermiticity of
\(A_n^{(2)}\), after multiplication by the invertible scalar denominators,
yields
\begin{align}
 &\bigl[q\{1-(2-q)r\}
 -(1-q)\{1-(1+q)r\}\bigr]\nonumber\\
 &\hspace{23mm}\times
 P_{q,r}Q_nP_{1-q,r}=0 .
\label{eq:bcgr-separator}
\end{align}
The coefficient factors exactly as
\((2q-1)(1-r)\).  It is nonzero for \(q\ne1/2\), because \(r\ne1\);
for \(q=1/2\), \(Q_1\) is scalar on the block.  Thus
\([Q_1,Q_n]=0\) for every \(n\), and the nonsingular CPOBC
central-generator lemma propagates commutativity to every transition
\cite{SrivastavaSurya2026}.

For the second statement, exact marked-isomorphism enumeration gives
\[
 8\ {\rm Markov}+26\ {\rm DGC}
 +35\ {\rm unequal\ CPOBC}+10\ {\rm equal\ CPOBC}.
\]
At a scalar reference the linearization \(J\) acts entrywise.  Its
three-dimensional kernel supplies
\(X_e=\sum_{\alpha=1}^3(n_\alpha)_eH_\alpha\).  Symmetric quadratic terms
are tangent to the three-parameter scalar solution manifold.  Modulo
\(\operatorname{im}J\), the remaining second-order equation is
\[
 \sum_{\alpha<\beta}
 [W_{\alpha\beta}]\otimes[H_\alpha,H_\beta]=0.
\]
The augmented rank in Eq.~\eqref{eq:bcgr-ranks} makes the three wedge classes
independent, so all three commutators vanish.  A certified nonzero
\(30\times30\) minor is
\[
 -\frac{1099511627776}{582076609134674072265625}.
\]
For the triangular ansatz, all diagonal equations hold by scalar
covariance.  Through four events a scalar CSG character uses the complete
three-jet \(\bm t=(t_1,t_2,t_3)\), with \(t_0=1\):
\[
 \lambda(v,m;\bm t)
 =\sum_{k=m}^{v}\binom{v-m}{k-m}t_k,\qquad
 a(c,P;\bm t)
 =\frac{\lambda(|P|,m(P);\bm t)}
 {\lambda(|c|,0;\bm t)}.
\]
Write the first three triangular antichain generators as
\[
 Q_i=\begin{pmatrix}a_i&x_i\\0&b_i\end{pmatrix}.
\]
An exact lower CPOBC residual has upper-right entry
\begin{align}
0={}&
\frac{(a_1a_2-a_1b_2-a_2)(a_2a_3-a_2b_3-a_3)}
{a_1a_2b_1b_2}\nonumber\\
&\times\big[(a_1-b_1)x_2-(a_2-b_2)x_1\big].
\label{eq:bcgr-triangular-separator}
\end{align}
For positive scalar transitions, both numerator factors are strictly
negative:
\[
a_1(a_2-b_2)-a_2<0,\qquad
a_2(a_3-b_3)-a_3<0.
\]
Thus \([Q_1,Q_2]=0\).  If the characters first differ at \(t_1\) or
\(t_2\), one global triangular similarity diagonalizes this lower pair,
and an exact \(22\times22\) late block has determinant
\begin{align}
D_1={}&
\frac{t_{L2}t_{R1}^{5}t_{R2}^{2}
(t_{L1}-t_{R1})(t_{R1}+t_{R2})^2}
{(1+t_{L1})(1+t_{R1})^{10}
(1+2t_{L1}+t_{L2})(1+2t_{R1}+t_{R2})^{12}},
\label{eq:bcgr-minor-t1}\\
D_2={}&
-\frac{t_1^{7}t_{R2}^{2}(t_1+t_{R2})^2(t_{L2}-t_{R2})}
{(1+t_1)^{10}(1+2t_1+t_{L2})
(1+2t_1+t_{R2})^{13}},
\label{eq:bcgr-minor-t2}
\end{align}
respectively.  The first expression is used when
\(t_{L1}\ne t_{R1}\); the second when \(t_{L1}=t_{R1}=t_1\) and
\(t_{L2}\ne t_{R2}\).  All undisplayed factors are positive on the
declared domain.

When the first difference is \(t_3\), the two characters agree through
stage two and the early extension coordinates cannot be removed by
similarity.  Exact block elimination instead gives
\begin{align}
D_{3,\mathrm{late}}
&=\frac{t_1^5t_2^2(t_1+t_2)^2}
{(1+t_1)^{10}(1+2t_1+t_2)^{12}},\\
D_{3,\mathrm{Schur}}
&=-\frac{t_1^3t_2(t_{L3}-t_{R3})^2}
{(1+t_1)^2(1+2t_1+t_2)}
\nonumber\\
&\hspace{7mm}\times
\frac{1}
{(1+3t_1+3t_2+t_{L3})^2
(1+3t_1+3t_2+t_{R3})^2}.
\label{eq:bcgr-minor-t3}
\end{align}
Their product is a nonzero \(29\times29\) minor.  The three
first-difference strata exhaust all distinct positive three-jets.  Since
the global-similarity coboundary is always present, the rank is exactly 29
and the kernel is exactly its span.

At five events there are 131 marked transitions and 820 residual
equations.  Suppose first that the characters differ among
\(t_1,t_2,t_3\).  The four-event restriction just proved is a global
coboundary.  Subtract it from all stages; the first 30 off-diagonal
coordinates vanish.  The remaining \(820\times101\) late matrix has exact
\(101\times101\) minors
\begin{align}
E_1={}&
\frac{t_{L2}t_{R1}^{37}t_{R2}^{13}t_{R3}^{3}
(t_{L1}-t_{R1})(t_{R1}+t_{R2})^{13}}
{(1+t_{L1})(1+t_{R1})^{44}
(1+2t_{L1}+t_{L2})(1+2t_{R1}+t_{R2})^{56}}
\nonumber\\[-1mm]
&\times
\frac{(t_{R2}+t_{R3})^8(t_{R1}+2t_{R2}+t_{R3})^6}
{(1+3t_{R1}+3t_{R2}+t_{R3})^{60}},
\label{eq:bcgr-five-t1}\\
E_2={}&-
\frac{t_1^{39}t_{R2}^{13}t_{R3}^{3}
(t_{L2}-t_{R2})(t_1+t_{R2})^{13}}
{(1+t_1)^{44}(1+2t_1+t_{L2})
(1+2t_1+t_{R2})^{57}}
\nonumber\\[-1mm]
&\times
\frac{(t_{R2}+t_{R3})^8(t_1+2t_{R2}+t_{R3})^6}
{(1+3t_1+3t_{R2}+t_{R3})^{60}},
\label{eq:bcgr-five-t2}\\
E_3={}&-
\frac{t_1^{38}t_2^{14}t_{R3}^{3}
(t_1+t_2)^{13}(t_2+t_{R3})^8}
{(1+t_1)^{44}(1+2t_1+t_2)^{56}}
\nonumber\\[-1mm]
&\times
\frac{(t_{L3}-t_{R3})(t_1+2t_2+t_{R3})^6}
{(1+3t_1+3t_2+t_{L3})
(1+3t_1+3t_2+t_{R3})^{61}} .
\label{eq:bcgr-five-t3}
\end{align}
The first formula applies when \(t_1\) first differs, the second when
\(t_1\) agrees and \(t_2\) first differs, and the third when \(t_1,t_2\)
agree and \(t_3\) first differs.  Every factor other than the declared
coupling difference is strictly positive.  These minors are independent of
\(t_{L4},t_{R4}\), so no stage-four coupling can restore a late extension.

It remains to treat a first difference in \(t_4\).  Put
\[
\Lambda_3=1+3t_1+3t_2+t_3,\qquad
\Lambda_{4\alpha}=1+4t_1+6t_2+4t_3+t_{\alpha4}.
\]
Omit one nonzero stage-four coboundary column.  A 100-dimensional late
block has determinant
\[
D_{4,\mathrm{late}}=
\frac{t_1^{37}t_2^{13}t_3^3(t_1+t_2)^{13}
(t_2+t_3)^8(t_1+2t_2+t_3)^6}
{(1+t_1)^{44}(1+2t_1+t_2)^{56}\Lambda_3^{60}}.
\]
Solving this block exactly and substituting into 30 remaining rows gives a
Schur determinant
\[
D_{4,\mathrm{Schur}}=
\frac{t_1^9t_2^4(t_1+t_2)^2(t_{L4}-t_{R4})^3}
{(1+t_1)^{12}(1+2t_1+t_2)^{13}
\Lambda_{4L}^3\Lambda_{4R}^3}.
\]
The first block is removed by 100 degree-one row pivots; 23 further leaf
pivots reduce the Schur block to a \(7\times7\) core.  Their product is a
nonzero \(130\times130\) minor,
\begin{equation}
\begin{split}
D_4={}&
\frac{t_1^{46}t_2^{17}t_3^3(t_1+t_2)^{15}
(t_2+t_3)^8(t_1+2t_2+t_3)^6}
{(1+t_1)^{56}(1+2t_1+t_2)^{69}\Lambda_3^{60}}\\
&\times
\frac{(t_{L4}-t_{R4})^3}
{\Lambda_{4L}^3\Lambda_{4R}^3}.
\end{split}
\label{eq:bcgr-five-t4}
\end{equation}
Thus every first-difference stratum has rank at least 130.  The global
similarity coboundary makes the rank at most 130, proving the third
statement through five events.

At six events the catalog contains
\[
2+6+22+101+576=707
\]
marked transitions and 9680 residual equations: 87 operator Markov sums,
4826 covariance equations, 3483 unequal-size CPOBC equations, and 1284
equal-size CPOBC equations.  If the first unequal coupling is
\(t_j\), \(j\le4\), the five-event restriction is already split.
Subtracting its global-similarity coboundary leaves the 576 stage-five
coordinates.  Exact \(576\times576\) late minors close all four strata.
The \(t_2\) minor is completely row-leaf eliminable; the \(t_1,t_3,t_4\)
minors have 573 degree-one column pivots and a \(3\times3\) core.  For
example, the \(t_3\) core is
\[
\begin{pmatrix}
-t_1/\Lambda_1&1/\Lambda_1&0\\
-t_2/\Lambda_2&0&1/\Lambda_2\\
0&-t_2/\Lambda_{3L}&t_1/\Lambda_{3R}
\end{pmatrix},
\]
where
\[
\Lambda_1=1+t_1,\quad
\Lambda_2=1+2t_1+t_2,\quad
\Lambda_{3\alpha}=1+3t_1+3t_2+t_{\alpha3}.
\]
Its determinant is proportional to \(t_{L3}-t_{R3}\); the \(t_1\) and
\(t_4\) cores have the same staircase form.  Multiplication by the exact
leaf factors gives
\[
F_j=(t_{Lj}-t_{Rj})\,P_j/Q_j,\qquad j=1,\ldots,4,
\]
with \(P_j,Q_j\) products of strictly positive couplings, consecutive
coupling sums, and CSG normalization polynomials.  Their complete
factorizations and row lists are retained in the supporting material
\cite{XuBellCausalGrowthRigidity2026}.

For the only remaining stratum, put
\[
t_{L1}=t_{R1}=t_1,\ldots,t_{L4}=t_{R4}=t_4,
\qquad t_{L5}\ne t_{R5},
\]
and define
\begin{align*}
\Lambda_1&=1+t_1,&\Lambda_2&=1+2t_1+t_2,\\
\Lambda_3&=1+3t_1+3t_2+t_3,\\
\Lambda_4&=1+4t_1+6t_2+4t_3+t_4,\\
\Lambda_{5\alpha}&=1+5t_1+10t_2+10t_3+5t_4+t_{\alpha5}.
\end{align*}
After one nonzero coboundary coordinate is omitted, 575 row leaves
eliminate the remaining late columns.  Four obstruction rows close the
four early tangent directions.  After scaling out
\((t_{L5}-t_{R5})/(\Lambda_{5L}\Lambda_{5R})\), their determinant is
\[
\frac{t_1^4t_2^2t_3}
{\Lambda_1\Lambda_2^3\Lambda_3^2\Lambda_4^2}>0.
\]
Restoring the four scaled factors and the late determinant gives the
nonzero \(706\times706\) minor
\begin{equation}
\begin{split}
D_5={}&-\frac{t_1^{304}t_2^{105}t_3^{26}t_4^4
(t_1+t_2)^{98}(t_2+t_3)^{71}(t_3+t_4)^{19}}
{\Lambda_1^{269}\Lambda_2^{309}\Lambda_3^{355}\Lambda_4^{361}}\\
&\times
\frac{(t_1+2t_2+t_3)^{49}(t_2+2t_3+t_4)^{40}}
{\Lambda_{5L}^4\Lambda_{5R}^4}\\
&\times
(t_1+3t_2+3t_3+t_4)^{22}(t_{L5}-t_{R5})^4 .
\end{split}
\label{eq:bcgr-six-t5}
\end{equation}
Every nondifference factor is positive.  Thus each of the five
first-difference strata has rank at least 706.  The coboundary supplies a
null vector, so the rank is exactly 706 and its span is the full kernel.

The six-event result is the first complete four-tangent audit of an
all-order antichain obstruction.  To derive the latter, let \(t_k\) be the
least unequal coupling and subtract the global-similarity coboundary so that
the upper-right entry of \(Q_k\) vanishes.  Below stage \(k\), the two
diagonal characters coincide and every triangular transition belongs to
the dual-number algebra
\(\mathbb C[\varepsilon]/(\varepsilon^2)\).  The nonsingular scalar CSG
classification uses only sums, products, and inverses of gregarious
transitions, so it base-changes to this algebra.  Consequently every early
coordinate is uniquely
\[
y_e=\sum_{\ell=1}^{k-1}c_\ell\partial_{t_\ell}a_e(\bm t).
\]

For \(k\ge3\), set
\(\Lambda_n=1+\sum_{j=1}^n\binom njt_j\).  The antichain pivots
\[
A_k^{(r)}A_r^{(0)}=A_r^{(r)}A_k^{(0)}
\]
solve the needed stage-\(k\) coordinates.  Substitute them into
\[
[A_k^{(1)},A_2^{(1)}]=0,\qquad
A_k^{(2)}A_2^{(1)}=A_2^{(2)}A_k^{(1)},
\]
and
\[
A_k^{(r)}A_{r+1}^{(1)}
=A_{r+1}^{(r)}A_k^{(1)},\qquad2\le r\le k-2.
\]
After scaling by
\(\Lambda_{kL}\Lambda_{kR}/(t_{Lk}-t_{Rk})\), one row interchange makes
the resulting \((k-1)\)-square matrix lower triangular, with diagonal
\[
-\frac{t_2}{\Lambda_1\Lambda_2},\quad
\frac{t_1^2}{\Lambda_2^2},\quad
\frac{t_1t_2}{\Lambda_3^2},\ldots,
\frac{t_1t_{k-2}}{\Lambda_{k-1}^2}.
\]
This proves Eq.~\eqref{eq:bcgr-allorder-det} and removes all early
tangents.  The cases \(k=1,2\) follow from
Eq.~\eqref{eq:bcgr-triangular-separator}.  The gregarious identity
\[
Q_nQ_1^{-1}Q_k=Q_kQ_1^{-1}Q_n\qquad(1<k<n)
\]
then diagonalizes every later antichain generator for \(k\ge2\); the
\(k=1\) case uses the complementary pair
\([Q_n,Q_2Q_1^{-1}]=0\) and
\([A_n^{(1)},A_2^{(1)}]=0\).  Atomization, Bell causality against the
gregarious alternative, and MSR propagate diagonality to all transitions.
Restoring the coboundary proves the all-order statement.

The symbolic archive derives the complete \(\mathcal O_k\) through
\(k=12\) and reproduces the six-event \(\mathcal O_5\) entry by entry.  A
second implementation using only exact rational arithmetic verifies
normalization cancellation, triangularity, Eq.~\eqref{eq:bcgr-allorder-det},
and positivity through \(k=40\).  These checks audit the formula; the
all-\(k\) conclusion follows from the displayed triangular factorization.

For the coincident-character branch, identify the upper-triangular algebra
with the dual numbers
\(\mathbb D=\mathbb C[\varepsilon]/(\varepsilon^2)\) through
\(\varepsilon\mapsto E_{12}\).  Strict positivity makes every denominator
in the scalar CSG induction a unit of \(\mathbb D\).  Because that induction
uses only sums, products, and inverses of gregarious transitions, it
base-changes to \(\mathbb D\): through parent size \(n\), the unique dual
couplings are
\[
 t_\ell+\varepsilon c_\ell,\qquad1\le\ell\le n.
\]
Restriction from \(n+1\) to \(n\) preserves the coefficients, giving one
infinite sequence \((c_\ell)_{\ell\ge1}\), and differentiation gives
Eq.~\eqref{eq:bcgr-equal-character}.  The image of \(\mathbb D\) is
commutative.  If \(c_k\) is the first nonzero coefficient, the nilpotent
entry of the gregarious transition from the \(k\)-antichain is
\(-c_k/\Lambda_k^2\ne0\), so the Jordan extension is nonsplit; the converse
is immediate.  An independent exact-rational implementation evaluates
47\,904 dual-number transition types through parent size 64.  The complete
four- and five-event systems have nullities three and four, exactly their
coupling-tangent dimensions, and all five tangent columns annihilate the
9680 six-event residuals.

For the irreducible two-dimensional branch, the gregarious identity makes
all \(R_n=Q_1^{-1}Q_n\) commute.  If \(R_2\) is a non-scalar Jordan block,
\([Q_1Q_2^{-1},Q_1^{-1}Q_2]=0\) makes \(Q_1\) preserve its unique
eigenline.  If \(R_2\) has simple spectrum, irreducibility leaves only
\[
 Q_1=\begin{pmatrix}0&1\\p&0\end{pmatrix},\qquad
 R_2=\operatorname{diag}(r,s),
\quad prs(r-s)(p-1)\ne0 .
\]
Writing \(R_3=\operatorname{diag}(u,v)\), two rows of the next CPOBC
relation have determinant
\begin{equation}
 -\frac{p^2(r-s)^2(r+s)[p(r^2+s^2)-rs]}
 {rs^2(p-1)}.
\label{eq:bcgr-early-irred}
\end{equation}
The first exceptional component \(s=-r\) forces \(u=v\).  At the following
stage the required two-precursor transition has determinant
\(-u^2D/r^2\), \(D=9pr^2-4r^2+1\), while two response rows have determinant
\(48pu^2/D\); nonsingularity therefore makes the latter full rank and
forces the next generator to vanish.  The other component has
\(p=rs/(r^2+s^2)\), \(v=-su/r\); away from the overlap \(r+s=0\), the
following two-row determinant is
\[
 \frac{2u^2(r-s)(r+s)^2}{r^3}\ne0.
\]
It gives the same contradiction.  A symbolic certificate reconstructs
these factorizations exactly.  This excludes non-scalar onset at \(R_2\),
not a delayed-onset branch with scalar \(R_2\).

For the surviving scalar ratio, \(Q_2=hQ_1\).  The exact transitions from
the two-element chain are
\begin{equation}
 G_2=hQ_1,\qquad V_2=h(\mathbf1-Q_1),\qquad
 T_2=(1-h)\mathbf1 .
\label{eq:bcgr-delayed-chain}
\end{equation}
Thus nonsingularity gives \(h\ne0,1\), eliminating every antichain
candidate on the \(h=1\) branch.  We next restrict
Eq.~\eqref{eq:bcgr-delayed-component} to the trace-balanced slice
\(d=a,\ x=-1\).  Direct substitution into all twelve
determinant-cleared stage-three residuals gives the common factor
\((y-1)(q-a^2)\).  On \(y\ne1\), nonsingularity permits removal of
\(q-a^2\), and two quotient rows give
\[
(2a-1)(a^2-a+q)^2=0,\qquad
(2a-1)^2(a^2-a+q)=0.
\]
The remaining rows force exactly \(E_1\) or \(E_2\) as stated in the
theorem; \(y=1\) is \(C_0\).  Direct factorization verifies sufficiency, so
these components exhaust the trace-balanced subchart over
\(\mathbb Q\).  The attempted exact elimination of \(d-a\) and \(x+1\)
from the full saturated first-simple-spectrum ideal did not produce a
finite characteristic-zero certificate; the complementary part of that
chart therefore remains open.

On \(C_0\), a putative next relative generator commutes with \(R_3\); write
\(R_5=\operatorname{diag}(1,z)\).  For three stage-four residuals
\(f_1,f_2,f_3\), with
\begin{equation}
 \Delta_3=-a^2h^2+6ah^2-6ah+h^2q-8h^2+18h-9 ,
\end{equation}
exact characteristic-zero reduction gives
\begin{equation}
 \operatorname{NF}_{I}\!\left[z(q-a^2)\Delta_3\right]\ne0,\qquad
 \operatorname{NF}_{I}\!\left([z(q-a^2)\Delta_3]^2\right)=0,
 \quad I=\langle f_1,f_2,f_3\rangle .
\label{eq:bcgr-delayed-radical}
\end{equation}
Hence every common zero has
\(z(q-a^2)\Delta_3=0\).  Yet \(z\ne0\),
\(q-a^2=-\det Q_1\ne0\), and
\(\Delta_3/h^2=\det A_3^{(3)}\ne0\) on the nonsingular branch, a
contradiction.  SymPy and Singular independently reproduce both normal
forms; saturation by all transition determinants gives the unit ideal.

On \(E_1,E_2\), a next-stage residual first forces \(y=\pm1\); \(y=1\)
returns to \(C_0\).  At \(y=-1\), two residuals force
\[
L_1=(2h-3)z-4h+3=0,\qquad
L_2=(3a-2)z-3a+1=0,
\]
respectively.  Their resultants with a third residual core are
\[
-(2h-3)(9h^2-29h+18),\qquad
(3a-2)(26a^2-31a+8).
\]
The linear leading factors cannot vanish with \(L_1,L_2\), while the
quadratics are required stage-three transition determinants.  Both
components are impossible.  Independent full saturated stage-four ideals
are the unit ideal over \(\mathbb Q\).

Transitive percolation is the subfamily \(t_k=t^k\).  With
\(p=t/(1+t)\), one compact corollary is
\[
 p_L^3p_R^{13}(p_R-p_L)^2(1-p_R)^{22}
 =
 \frac{t_L^3t_R^{13}(t_R-t_L)^2}
 {(1+t_L)^5(1+t_R)^{37}} .
\]
At the archived rational check pair
\((t_L,t_R)=(1/4,1/2)\), this minor is
\[
 \frac{16777216}{1407137205909366759375}.
\]
\end{proof}

Equation~\eqref{eq:bcgr-hermiticity} alone cannot prove the first part.  For
\(Q_1=\operatorname{diag}(\lambda,1-\lambda)\), one positive off-diagonal
\(Q_n\) can make Eq.~\eqref{eq:bcgr-one-precursor} positive while
\([Q_1,Q_n]\ne0\).  This local example is a failure control for the
one-precursor argument, not a complete CPOBC growth law: the second
antichain-generator identity, nonsingularity of the chain timid transition,
and the two-precursor self-adjointness condition supply the missing
separators.  The second part excludes a noncommuting
\emph{leading} tangent through second order; it does not exclude
higher-order contact or a disconnected irreducible component in the
non-self-adjoint or singular branches.  The third part applies to every
pair of distinct strictly positive scalar CSG coupling sequences,
irrespective of the stage at which they first differ.  It does not cover
zero/singular coupling boundaries, irreducible nontriangular
representations, or higher-dimensional noncommutative nilpotent structure.
The fourth part closes the positive equal-character \(2\times2\)
upper-triangular branch.  The fifth closes non-scalar onset in an
irreducible nonsingular \(2\times2\) branch.  The sixth removes the scalar
value \(h=1\) and exhausts the trace-balanced \(d=a,\ x=-1\) subchart at
\(R_3\), but does not classify the rest of the first-simple-spectrum chart,
Jordan onset, or first onset after \(R_3\).

This proves A11BR as a boundary theorem.  It narrows the square-operator
search to non-self-adjoint nonsingular transitions, zero/singular
couplings, state-only covariance, higher-dimensional noncommutative
nilpotent thickenings, the unclassified part of the normalized
first-simple-spectrum chart, and the remaining delayed-onset irreducible
nontriangular components with scalar \(R_2=h\mathbf1\), \(h\ne0,1\),
including Jordan onset and later first onset.
It neither constructs nor excludes such a representation.  Nor does it
provide an infinite noncommutative CPOBC measure with consistent cylinder
marginals, a Lorentzian continuum, or geometry backreaction.

\subsection{Falsifying deformations and the remaining assumption}

Three controls isolate the mechanism.  With only one internal mode,
\(\Str M^2=-\Delta\), so the logarithmic cutoff dependence remains.
Erasing parity replaces the supertrace by a positive trace with
\(\Tr\mathbf1=4\), restoring the quadratic divergence.  Most importantly,
the lowest number-conserving nonadditive deformation,
\begin{equation}
 M_\eta^2=x+\Delta N+\eta N(N-1),
 \label{eq:ph-deformation}
\end{equation}
gives
\begin{equation}
 \Str M_\eta^2=2\eta .
 \label{eq:ph-deformation-moment}
\end{equation}
The free result is therefore structural but not generically
interaction-stable.  A partial Ward identity or renormalization-group
invariant manifold must forbid or compensate the \(N(N-1)\) term without
enforcing the complete pairing ruled out by
Theorem~\ref{thm:ph-pairing-obstruction}.

This distinction separates eight statements.  A11BM is proved by the
positive finite-graph multiplet.  A11BW is proved by
Theorem~\ref{thm:selective-ladder-closure} and the finite geometry channel:
the required moments are selectively protected for arbitrary finite
positive interacting \(A_\omega\), with a quantitative defect bound, and
the resulting weights have a local positive dynamics.  A11PG is proved by
Theorem~\ref{thm:pcg-projective-law}: on a supplied naturally labelled
causal-birth language, those weights generate compatible finite cylinder
laws and an infinite positive history measure.  A11CG is proved by
Theorems~\ref{thm:cqg-positive-growth}
and~\ref{thm:cqg-complex-extension}: the response-generated contact
sequence is DGC and Bell causal, and one \(r\geq2\) coherent subfamily has a
countably additive strongly positive rank-one quantum measure.  A11NH is
proved by Theorem~\ref{thm:nonabelian-history-extension}: on a supplied
two-port strip, a fresh-ancilla local unitary generates an extendible
rank-four non-Abelian quantum measure.  A11CH is proved by
Theorem~\ref{thm:covariant-nonabelian-holonomy}: on arbitrary CSG causal
histories, intrinsic fresh-event gates define a covariant countably
additive non-Abelian operator measure and strongly positive decoherence
functional.  A11OC is proved by
Theorem~\ref{thm:central-operator-growth-classification}: central
operator-valued CSG has an infinite POVM and exact covariance, but is
classified as a coherence-blind direct integral of scalar growth laws; the
positive noncentral control proves that one-step normalization alone does
not preserve DGC.  A11BR is proved by
Theorem~\ref{thm:bcgr-rigidity}: every finite-dimensional self-adjoint
nonsingular CPOBC system classicalizes, without a sign assumption and
including complementary-spectrum blocks; independently, the tested scalar
locus has no noncommuting leading tangent, and every \(2\times2\)
upper-triangular extension between distinct strictly positive scalar CSG
coupling sequences splits, at arbitrary first-difference order, while every
coincident-character \(2\times2\) self-extension is a commuting coupling
tangent; an irreducible nonsingular \(2\times2\) branch must have scalar
\(R_2=h\mathbf1\), with \(h\ne0,1\), and the trace-balanced
\(d=a,\ x=-1\) subchart of the normalized first-simple-spectrum chart at
\(R_3\) is exhausted and obstructed at the next antichain stage.
These are exact exclusions, not an existence theorem or a no-go
theorem for non-self-adjoint, singular, higher-dimensional noncommutative
nilpotent, the remainder of the first-simple-spectrum chart, the remaining
delayed-onset irreducible nontriangular, or state-only formulations.
A11B remains open and is stronger:
the \emph{same interacting Lorentzian
dynamics} must generate the relational records, growth language, topology,
contact response, and arity, replace the external classical CSG density by
a genuinely quantum backreacting geometry law satisfying the desired
operator locality/constraint conditions; derive rather than declare the
protected commutant; produce a
tight Lorentzian continuum measure; satisfy spin--statistics and anomaly
freedom; and couple consistently to nonlinear gauge constraints and graph
backreaction.

\subsection{Operational completion criterion for nonperturbative gravity}
\label{sec:nonperturbative-completion-criterion}

To prevent a growing list of model closures from being mistaken for a
nonperturbative definition, we record the minimum completion target used by
the present program.  A refining family \(\{\mathfrak Q_N\}\) will count as
a nonperturbative RQCP-QG completion only if one construction supplies all
of the following, with compatible orders of limits:
\begin{enumerate}[label=(NP\arabic*)]
\item a normalized positive microscopic process or state on each finite
  regulator, with relational rather than external time;
\item autonomous fluctuating adjacency, causal order, additive physical
  volume, and a tight nondegenerate probability law on the resulting
  geometries;
\item a projectively consistent continuum net of local observables with a
  positive Lorentzian reconstruction or an equivalent constructive
  continuation;
\item anomaly-free nonlinear diffeomorphism constraints, expressed by a
  closed BV/BRST master equation or an equivalent exact relational Ward
  algebra;
\item regulator-independent relational observables, including matter
  correlators and the Newton response, with graph topology changes,
  backreaction, and caustic sectors controlled rather than excluded;
\item a semiclassical regime in which the same observables recover
  nonlinear Einstein--matter dynamics with an error that vanishes
  uniformly on a stated physical domain.
\end{enumerate}
Every item is a simultaneous requirement, not a menu.  The present series
has explicit finite or quadratic model realizations of parts of (NP1),
(NP2), (NP4), and (NP6), and the regulator theorem plus the present
multiplet and selective-Ward channel close one finite interacting,
fixed-language contribution to (NP5).  The projective-growth channel closes
the positive finite-to-infinite probability part of (NP2), and the
covariant-growth companion removes the natural-label/Bell defect and
supplies one extendible rank-one coherent measure.  The fixed-strip
companion separately supplies an extendible local-unitary non-Abelian
measure.  The causal-holonomy companion combines arbitrary-causet CSG
covariance with an intrinsic noncommuting operator density and proves a
countably additive measure whose integrated range remains non-Abelian.
The central operator-growth classification supplies an infinite geometry
POVM but proves that its probabilities are only direct mixtures; its
positive noncentral control fails DGC.  The CPOBC rigidity theorem further
excludes the complete finite-dimensional self-adjoint nonsingular branch
without a sign assumption, every noncommuting leading tangent at one exact
four-event scalar point, and every nonsplit \(2\times2\)
upper-triangular extension between distinct strictly positive scalar CSG
coupling sequences; it also classifies every coincident-character
\(2\times2\) self-extension as a commuting coupling tangent.  It leaves
non-self-adjoint nonsingular, zero/singular-coupling, state-only,
higher-dimensional noncommutative nilpotent, delayed-onset irreducible
nontriangular components with scalar
\(R_2=h\mathbf1\), \(h\ne0,1\), Jordan/later onset, and infinite branches
open.
These results do not autonomously
select the language, contact sequence, or arity, produce a genuinely
noncentral operator-Markov/Bell quantum-backreacting law, or provide
continuum tightness.
The series has no single
construction satisfying (NP1)--(NP6), no projectively tight interacting
Lorentzian continuum measure, and no nonlinear anomaly-free continuum
constraint algebra.  Accordingly these bridges shorten the open cut but do
not constitute nonperturbative quantum gravity.
